\documentclass[10pt,aspectratio=169]{beamer}

% ============================================================
% Módulo 3: Construcción de Escenarios Macro-Financieros
% Stress Testing en Sistemas de Pensiones de Capitalización Individual
% MACROPOL
% ============================================================

\usetheme{metropolis}
\usecolortheme{seahorse}
\usefonttheme{professionalfonts}

\usepackage[spanish,es-noquoting]{babel}
\usepackage[utf8]{inputenc}
\usepackage[T1]{fontenc}
\usepackage{amsmath,amssymb,bm}
\usepackage{booktabs}
\usepackage{array}
\usepackage{graphicx}
\usepackage{tikz}
\usepackage{listings}
\usepackage{xcolor}
\usetikzlibrary{arrows.meta,positioning,calc}

% Colores MACROPOL
\definecolor{MacropolBlue}{RGB}{0,67,122}
\definecolor{MacropolTeal}{RGB}{0,134,150}
\definecolor{MacropolGray}{RGB}{80,80,80}
\definecolor{SoftBlue}{RGB}{235,244,250}
\definecolor{SoftRed}{RGB}{252,238,238}
\definecolor{SoftGreen}{RGB}{235,248,240}
\definecolor{CodeBg}{RGB}{248,248,248}

\setbeamercolor{title}{fg=MacropolBlue}
\setbeamercolor{frametitle}{fg=MacropolBlue}
\setbeamercolor{structure}{fg=MacropolBlue}
\setbeamercolor{block title}{bg=MacropolBlue,fg=white}
\setbeamercolor{block body}{bg=SoftBlue,fg=black}
\setbeamertemplate{navigation symbols}{}
\setbeamertemplate{footline}[frame number]

% R code style
\lstdefinelanguage{R}{
  keywords={function,if,else,for,while,repeat,in,next,break,TRUE,FALSE,NULL,NA,NaN,Inf,library},
  otherkeywords={<-,<<-,$,~},
  sensitive=true,
  morecomment=[l]\#,
  morestring=[b]",
  morestring=[b]'
}

\lstset{
    language=R,
    basicstyle=\ttfamily\scriptsize,
    keywordstyle=\color{MacropolBlue}\bfseries,
    commentstyle=\color{MacropolGray}\itshape,
    stringstyle=\color{MacropolTeal},
    backgroundcolor=\color{CodeBg},
    frame=single,
    rulecolor=\color{MacropolGray!40},
    breaklines=true,
    showstringspaces=false,
    columns=fullflexible
}

% Macros
\newcommand{\baseline}{\textcolor{MacropolTeal}{\textbf{baseline}}}
\newcommand{\adverso}{\textcolor{red!70!black}{\textbf{adverso}}}
\newcommand{\R}{\texttt{R}}
\newcommand{\dd}{\mathrm{d}}

% Title
\title[Stress Testing Pensiones -- Módulo 3]{Módulo 3: Construcción de Escenarios Macro-Financieros}
\subtitle{ARIMA, VAR, Monte Carlo y diseño de choques para sistemas de pensiones}
\author{MACROPOL: Economía Aplicada y Políticas Públicas}
\institute{Stress Testing en Sistemas de Pensiones de Capitalización Individual}
\date{Abril 2026}

\begin{document}

% ============================================================
\begin{frame}
\titlepage
\end{frame}

% ============================================================
\begin{frame}{Dónde estamos en el curso}
\begin{block}{Módulo 3 dentro del flujo de stress testing previsional}
Transformamos una \textbf{narrativa macroeconómica} en trayectorias cuantitativas de variables que alimentan el stress test del portafolio, el módulo actuarial y el reporte regulatorio.
\end{block}

\begin{center}
\begin{tikzpicture}[node distance=0.5cm, every node/.style={font=\scriptsize}]
\node[draw, rounded corners, fill=SoftBlue, minimum width=2.4cm, minimum height=0.8cm] (m1) {M1 Fundamentos};
\node[draw, rounded corners, fill=SoftBlue, right=of m1, minimum width=2.6cm, minimum height=0.8cm] (m2) {M2 Riesgos};
\node[draw, rounded corners, fill=SoftGreen, right=of m2, minimum width=3.0cm, minimum height=0.8cm] (m3) {\textbf{M3 Escenarios}};
\node[draw, rounded corners, fill=SoftBlue, right=of m3, minimum width=2.7cm, minimum height=0.8cm] (m4) {M4 Portafolio};
\node[draw, rounded corners, fill=SoftBlue, below=0.7cm of m4, minimum width=2.7cm, minimum height=0.8cm] (m5) {M5 Actuarial};
\node[draw, rounded corners, fill=SoftBlue, left=of m5, minimum width=3.0cm, minimum height=0.8cm] (m6) {M6 Reporte};

\draw[-{Latex}, thick] (m1) -- (m2);
\draw[-{Latex}, thick] (m2) -- (m3);
\draw[-{Latex}, thick] (m3) -- (m4);
\draw[-{Latex}, thick] (m4) -- (m5);
\draw[-{Latex}, thick] (m5) -- (m6);
\end{tikzpicture}
\end{center}

\textbf{Producto del módulo:} matriz de escenarios macro-financieros por horizonte, variable y severidad.
\end{frame}

% ============================================================
\begin{frame}{Objetivos de aprendizaje}
Al finalizar el módulo, los participantes podrán:

\begin{enumerate}
  \item Construir un escenario \baseline{} y uno \adverso{} para variables macro-financieras relevantes.
  \item Estimar modelos ARIMA para proyecciones univariadas.
  \item Estimar modelos VAR para capturar interdependencias macro-financieras.
  \item Generar trayectorias Monte Carlo y bandas de incertidumbre.
  \item Calibrar choques históricos e hipotéticos con trazabilidad regulatoria.
  \item Implementar un flujo reproducible en \R{}.
\end{enumerate}

\vspace{0.4em}
\begin{block}{Principio conductor}
No buscamos ``predecir el futuro''. Buscamos construir trayectorias \textbf{coherentes, severas y plausibles} para evaluar la resiliencia del sistema previsional.
\end{block}
\end{frame}

% ============================================================
\begin{frame}{Ruta gradual del módulo}
\begin{center}
\begin{tikzpicture}[node distance=0.35cm, every node/.style={font=\scriptsize}]
\node[draw, rounded corners, fill=SoftBlue, minimum width=2.5cm, minimum height=0.8cm] (n1) {1. Datos};
\node[draw, rounded corners, fill=SoftBlue, right=of n1, minimum width=2.5cm, minimum height=0.8cm] (n2) {2. ARIMA};
\node[draw, rounded corners, fill=SoftBlue, right=of n2, minimum width=2.5cm, minimum height=0.8cm] (n3) {3. VAR};
\node[draw, rounded corners, fill=SoftBlue, right=of n3, minimum width=2.5cm, minimum height=0.8cm] (n4) {4. Monte Carlo};
\node[draw, rounded corners, fill=SoftBlue, below=0.55cm of n4, minimum width=2.5cm, minimum height=0.8cm] (n5) {5. Choques};
\node[draw, rounded corners, fill=SoftGreen, left=of n5, minimum width=2.5cm, minimum height=0.8cm] (n6) {6. Taller R};

\draw[-{Latex}, thick] (n1) -- (n2);
\draw[-{Latex}, thick] (n2) -- (n3);
\draw[-{Latex}, thick] (n3) -- (n4);
\draw[-{Latex}, thick] (n4) -- (n5);
\draw[-{Latex}, thick] (n5) -- (n6);
\end{tikzpicture}
\end{center}

\begin{columns}
\column{0.48\textwidth}
\begin{block}{De menor a mayor complejidad}
\begin{itemize}
  \item Una variable: ARIMA.
  \item Varias variables: VAR.
  \item Muchas trayectorias: Monte Carlo.
  \item Narrativa adversa: choques calibrados.
\end{itemize}
\end{block}
\column{0.48\textwidth}
\begin{block}{De modelo a decisión}
\begin{itemize}
  \item Escenarios reproducibles.
  \item Supuestos auditables.
  \item Impactos comparables.
  \item Reglas de acción supervisora.
\end{itemize}
\end{block}
\end{columns}
\end{frame}

% ============================================================
\section{1. Contexto}

% ============================================================
\begin{frame}{¿Qué es un escenario macro-financiero?}
\begin{block}{Definición operativa}
Un escenario es una trayectoria conjunta de variables macroeconómicas y financieras bajo una narrativa específica, con horizonte, frecuencia, severidad y supuestos explícitos.
\end{block}

\begin{table}
\small
\begin{tabular}{p{3.2cm}p{4.2cm}p{4.2cm}}
\toprule
Componente & \baseline{} & \adverso{} \\
\midrule
Narrativa & Normalización gradual & Desaceleración + tasas altas + depreciación \\
Modelo & Proyección central & Proyección condicionada a choques \\
Supuestos & Política y mercado sin disrupciones & Persistencia de inflación y prima de riesgo \\
Resultado & Trayectoria esperada & Pérdida de valor y deterioro de aportes \\
\bottomrule
\end{tabular}
\end{table}
\end{frame}

% ============================================================
\begin{frame}{Variables relevantes para pensiones de capitalización individual}
\begin{columns}
\column{0.52\textwidth}
\begin{block}{Canal financiero}
\begin{itemize}
  \item Tasa de interés local.
  \item Inflación y tasa real.
  \item Tipo de cambio.
  \item Spread soberano.
  \item Rentabilidad de bonos y renta variable.
\end{itemize}
\end{block}

\column{0.42\textwidth}
\begin{block}{Canal previsional}
\begin{itemize}
  \item Salario real.
  \item Empleo formal.
  \item Densidad de cotización.
  \item Aportes.
  \item Saldo acumulado.
\end{itemize}
\end{block}
\end{columns}

\vspace{0.4em}
\begin{center}
\begin{tikzpicture}[node distance=0.6cm, every node/.style={font=\scriptsize}]
\node[draw, rounded corners, fill=SoftRed, minimum width=3cm, minimum height=0.8cm] (shock) {Choque macro};
\node[draw, rounded corners, fill=SoftBlue, right=of shock, minimum width=3.4cm, minimum height=0.8cm] (risk) {Factores de riesgo};
\node[draw, rounded corners, fill=SoftGreen, right=of risk, minimum width=3.2cm, minimum height=0.8cm] (pension) {Balance previsional};
\draw[-{Latex}, thick] (shock) -- (risk);
\draw[-{Latex}, thick] (risk) -- (pension);
\end{tikzpicture}
\end{center}
\end{frame}

% ============================================================
\section{2. ARIMA}

% ============================================================
\begin{frame}{ARIMA: una variable, su historia}
\begin{block}{¿Cuándo usar ARIMA en escenarios previsionales?}
ARIMA es el punto de entrada para construir trayectorias \baseline{} de una variable individual cuando queremos capturar persistencia, reversión y shocks idiosincráticos.
\end{block}

Modelo ARIMA$(p,d,q)$:
\[
\phi(L)(1-L)^d y_t = c + \theta(L)\varepsilon_t, \quad \varepsilon_t \sim (0,\sigma^2)
\]

\begin{block}{Intuición económica}
El pasado reciente contiene información sobre la velocidad de reversión, persistencia inflacionaria o inercia de tasas. ARIMA formaliza esa memoria.
\end{block}

\small \textit{Limitación:} No impone coherencia entre variables. Por eso usamos VAR después.
\end{frame}

% ============================================================
\begin{frame}[fragile]{Demo 1: ARIMA en R}
\begin{lstlisting}
library(forecast)
library(ggplot2)

# Serie trimestral de inflacion (simulada)
set.seed(42)
inflacion <- rnorm(40, mean = 4.5, sd = 1.2)
infl_ts <- ts(inflacion, start = c(2016, 1), frequency = 4)

# Seleccion automatica
fit_arima <- auto.arima(infl_ts, seasonal = FALSE)
summary(fit_arima)

# Pronostico baseline 8 trimestres
fc_base <- forecast(fit_arima, h = 8, level = c(50, 80, 95))
autoplot(fc_base)  # baseline + bandas de incertidumbre

# Crear escenario adverso: baseline + 1.5*sigma
fc_adverso <- fc_base$mean + 1.5 * sd(residuals(fit_arima))
\end{lstlisting}

\textbf{Nota del instructor:} Este es tu primer escenario base. Verás que ARIMA genera bandas de incertidumbre. Luego sumaremos shocks para el adverso.
\end{frame}

% ============================================================
\section{3. VAR}

% ============================================================
\begin{frame}{VAR: múltiples variables coherentes}
\begin{block}{¿Por qué pasar de ARIMA a VAR?}
En pensiones, los shocks no llegan aislados: inflación, tasas, tipo de cambio, PIB y spreads se mueven conjuntamente. VAR captura esa propagación empírica.
\end{block}

Sistema VAR$(p)$:
\[
\bm y_t = \bm c + A_1 \bm y_{t-1} + \cdots + A_p \bm y_{t-p} + \bm u_t, \quad \bm u_t \sim (0,\Sigma_u)
\]

\begin{block}{Ventajas}
\begin{itemize}
  \item Coherencia entre variables.
  \item Propagación dinámica de shocks.
  \item Impulse Response Functions (IRF) para leer canales.
  \item Útil para sensibilidad supervisora.
\end{itemize}
\end{block}
\end{frame}

% ============================================================
\begin{frame}[fragile]{Demo 2: VAR en R}
\begin{lstlisting}
library(vars)
library(tidyverse)

# Generar 3 variables correlacionadas: inflacion, tasa, fx
set.seed(42)
n <- 40
rho <- matrix(c(1.0, 0.3, 0.5,
                0.3, 1.0, 0.7,
                0.5, 0.7, 1.0), 3, 3)
Z <- matrix(rnorm(n * 3), n, 3)
L <- chol(rho)
datos <- Z %*% t(L)
colnames(datos) <- c("inflacion", "tasa", "fx")

# Estimar VAR(2)
fit_var <- VAR(datos, p = 2, type = "const")
summary(fit_var)

# Impulse response: efecto de shock en FX sobre inflacion y tasa
irf_fx <- irf(fit_var, impulse = "fx", response = c("inflacion", "tasa"),
              n.ahead = 8, boot = TRUE, ci = 0.90)
plot(irf_fx)

# Pronostico baseline multivariado
fc_var <- predict(fit_var, n.ahead = 8, ci = 0.95)
\end{lstlisting}

\textbf{Nota del instructor:} Observa cómo el VAR propaga automáticamente el shock de FX a las otras variables según sus dinámicas estimadas.
\end{frame}

% ============================================================
\section{4. Monte Carlo}

% ============================================================
\begin{frame}{Monte Carlo: de una trayectoria a muchas}
\begin{block}{¿Por qué Monte Carlo?}
Un solo escenario oculta incertidumbre. Monte Carlo permite generar miles de trayectorias compatibles con el modelo y medir distribución de resultados, percentiles y probabilidades.
\end{block}

Simulamos shocks aleatorios $\bm u_t^{(m)} \sim \mathcal{N}(0,\hat \Sigma_u)$ para $m = 1,\ldots,M$ (típicamente 5,000--10,000):

\[
\bm y_{t+h}^{(m)} = \hat{\bm c} + \hat A_1 \bm y_{t+h-1}^{(m)} + \cdots + \hat A_p \bm y_{t+h-p}^{(m)} + \bm u_{t+h}^{(m)}
\]

\begin{block}{Resultado}
Distribución de valores finales del fondo por percentil (5\%, 50\%, 95\%), no solo media.
\end{block}
\end{frame}

% ============================================================
\begin{frame}[fragile]{Demo 3: Monte Carlo en R}
\begin{lstlisting}
library(MASS)

# Parametros del VAR
K <- ncol(datos)
H <- 8
M <- 5000
Sigma <- summary(fit_var)$covres

# Simulacion Monte Carlo simple
sim_paths <- array(NA, dim = c(H, K, M))

for (m in 1:M) {
  y_lag <- tail(datos, 2)
  for (h in 1:H) {
    shock <- mvrnorm(1, mu = rep(0, K), Sigma = Sigma)
    # Aqui entra la formula del VAR completa (implementada en taller)
    y_new <- colMeans(datos) + shock
    sim_paths[h, , m] <- y_new
  }
}

# Percentiles por horizonte y variable
p05 <- apply(sim_paths, c(1, 2), quantile, probs = 0.05)
p50 <- apply(sim_paths, c(1, 2), quantile, probs = 0.50)
p95 <- apply(sim_paths, c(1, 2), quantile, probs = 0.95)

# Visualizar: trayectoria mediana vs colas
matplot(p50, type = "l", lty = 1, col = 1:K)
matlines(p05, lty = 2, col = 1:K); matlines(p95, lty = 2, col = 1:K)
\end{lstlisting}

\textbf{Nota del instructor:} Observa cómo los percentiles divergen: la incertidumbre crece con el horizonte.
\end{frame}

% ============================================================
\section{5. Choques}

% ============================================================
\begin{frame}{Choques: históricos vs. hipotéticos}
\begin{table}
\small
\begin{tabular}{p{3.0cm}p{4.2cm}p{4.2cm}}
\toprule
Tipo & Ventaja & Riesgo \\
\midrule
Histórico & Observable, comunicable, trazable & Puede no repetir estructura actual \\
Hipotético & Diseñado para vulnerabilidades actuales & Puede parecer arbitrario \\
Híbrido & Combina severidad observada y narrativa actual & Requiere mayor documentación \\
\bottomrule
\end{tabular}
\end{table}

\begin{block}{Principio}
El shock debe ser suficientemente severo para revelar vulnerabilidades, pero suficientemente plausible para guiar decisiones.
\end{block}

\textbf{Dos escenarios adversos que construiremos hoy:}
\begin{enumerate}
  \item Escenario A: Calibrado con crisis 2008 (histórico).
  \item Escenario B: Diseñado por narrativa de tasas altas + depreciación (hipotético).
\end{enumerate}
\end{frame}

% ============================================================
\begin{frame}[fragile]{Demo 4: Escenarios adversos}
\begin{lstlisting}
# ESCENARIO BASE: percentil 50 de Monte Carlo
escenario_base <- p50

# ESCENARIO ADVERSO HISTORICO (Crisis 2008)
# Caida: inflacion +1.5sd, tasa +150pb, fx +15%
shock_historico <- c(
  inflacion = +1.5,    # puntos porcentuales
  tasa      = +1.5,
  fx        = +15
)
escenario_hist <- t(t(p50) + shock_historico)

# ESCENARIO ADVERSO HIPOTETICO (Tasas + depreciacion)
# Narrativa: endurecimiento monetario y salida de capitales
shock_hip <- c(
  inflacion = +2.0,    # inflacion mas persistente
  tasa      = +2.5,    # suba mas fuerte
  fx        = +20      # depreciacion severa
)
escenario_hip <- t(t(p50) + shock_hip)

# Comparar: impacto en rentabilidad real
# r_real = tasa - inflacion
r_real_base <- escenario_base[, "tasa"] - escenario_base[, "inflacion"]
r_real_hist <- escenario_hist[, "tasa"] - escenario_hist[, "inflacion"]
r_real_hip  <- escenario_hip[, "tasa"] - escenario_hip[, "inflacion"]

data.frame(horizonte = 1:H,
           base = r_real_base,
           historico = r_real_hist,
           hipotetico = r_real_hip)
\end{lstlisting}
\end{frame}

% ============================================================
\section{6. Taller}

% ============================================================
\begin{frame}{Taller práctico: producto final}
\begin{block}{Entregable}
Un script reproducible en \R{} que genere una \textbf{matriz de escenarios macro-financieros} lista para Módulo 4.
\end{block}

\begin{table}
\tiny
\begin{tabular}{lllll}
\toprule
horizonte & escenario & inflacion & tasa & fx \\
\midrule
1 & baseline & 4.5 & 6.0 & 3.0 \\
1 & adverso\_historico & 6.0 & 7.5 & 18.0 \\
1 & adverso\_hipotetico & 6.5 & 8.5 & 23.0 \\
2 & baseline & 4.6 & 6.1 & 3.1 \\
\vdots & \vdots & \vdots & \vdots & \vdots \\
\bottomrule
\end{tabular}
\end{table}

\begin{block}{Estructura sugerida}
\texttt{00\_paquetes.R} → \texttt{01\_datos.R} → \texttt{02\_arima.R} → \texttt{03\_var.R} → \texttt{04\_montecarlo.R} → \texttt{05\_escenarios.R} → \texttt{06\_exportar.csv}
\end{block}
\end{frame}

% ============================================================
\begin{frame}{Guía de trabajo en grupos}
\begin{table}
\small
\begin{tabular}{p{2.7cm}p{4.2cm}p{4.2cm}}
\toprule
Grupo & Foco técnico & Pregunta de decisión \\
\midrule
A & ARIMA inflación y tasa & ¿El baseline es razonable? \\
B & VAR macro-financiero & ¿El shock se propaga coherentemente? \\
C & Monte Carlo & ¿Qué percentil debe usarse como adverso? \\
D & Choque histórico/hipotético & ¿Cuál escenario es más relevante? \\
\bottomrule
\end{tabular}
\end{table}

\begin{block}{Discusión plenaria}
Cada grupo presenta una \textbf{recomendación regulatoria}, no solo un resultado estadístico.
\end{block}
\end{frame}

% ============================================================
\begin{frame}{Checklist regulatorio del escenario}
\begin{enumerate}
  \item ¿Está claramente separado el \baseline{} del \adverso{}?
  \item ¿La narrativa económica explica el origen del shock?
  \item ¿El modelo cuantitativo es replicable?
  \item ¿Los supuestos están documentados?
  \item ¿La severidad es extrema pero plausible?
  \item ¿Las variables son relevantes para pensiones?
  \item ¿El horizonte coincide con el uso del stress test?
  \item ¿Los resultados se exportan para portafolio y actuarial?
  \item ¿Existe análisis de sensibilidad?
  \item ¿La conclusión guía una decisión regulatoria?
\end{enumerate}
\end{frame}

% ============================================================
\begin{frame}{Cierre: mapa de decisión}
\begin{center}
\begin{tikzpicture}[node distance=0.65cm, every node/.style={font=\scriptsize}]
\node[draw, rounded corners, fill=SoftBlue, minimum width=3.0cm, minimum height=0.8cm] (arima) {ARIMA};
\node[draw, rounded corners, fill=SoftBlue, right=of arima, minimum width=3.0cm, minimum height=0.8cm] (var) {VAR};
\node[draw, rounded corners, fill=SoftBlue, right=of var, minimum width=3.0cm, minimum height=0.8cm] (mc) {Monte Carlo};
\node[draw, rounded corners, fill=SoftRed, below=0.85cm of var, minimum width=3.3cm, minimum height=0.8cm] (shock) {Choques};
\node[draw, rounded corners, fill=SoftGreen, below=0.85cm of shock, minimum width=4.0cm, minimum height=0.8cm] (esc) {Escenario macro-financiero};
\node[draw, rounded corners, fill=SoftBlue, below=0.85cm of esc, minimum width=4.0cm, minimum height=0.8cm] (decision) {Decisión regulatoria};

\draw[-{Latex}, thick] (arima) -- (var);
\draw[-{Latex}, thick] (var) -- (mc);
\draw[-{Latex}, thick] (var) -- (shock);
\draw[-{Latex}, thick] (mc) -- (shock);
\draw[-{Latex}, thick] (shock) -- (esc);
\draw[-{Latex}, thick] (esc) -- (decision);
\end{tikzpicture}
\end{center}

\begin{block}{Mensaje final}
El objetivo del módulo no es elegir el modelo ``más complejo'', sino construir escenarios que sean \textbf{defendibles, reproducibles y útiles} para proteger la suficiencia previsional.
\end{block}
\end{frame}

% ============================================================
\end{document}